\section{Diseño de la investigación}

En esta sección se establece el marco metodológico que rige el desarrollo del presente trabajo, definiendo la 
tipología de la investigación, el enfoque de ingeniería adoptado para el desarrollo del software y la estrategia 
de validación del sistema.

\subsection{Naturaleza y tipo de investigación}
El presente trabajo se clasifica como una \textbf{investigación aplicada con enfoque tecnológico y desarrollo 
experimental}. Es aplicada porque busca resolver un problema concreto en el área de la Inteligencia Artificial 
Explicable (XAI): la limitación de la representación bidimensional para la inspección de capas convolucionales. 
Asimismo, es de carácter tecnológico y experimental, ya que el resultado principal consiste en la formulación 
de una metodología y la construcción de una herramienta de software funcional (\textit{artifact}) mediante la 
integración de Python, PyTorch, OpenVDB y Blender.

\subsection{Metodología de desarrollo de software}
Para la concepción, diseño e implementación del complemento y el pipeline de extracción, se adoptó una 
adaptación del ciclo de vida iterativo basado en el \textbf{Diseño Basado en Componentes} y el modelo de 
prototipado rápido:

\begin{enumerate}
    \item \textbf{Fase de Análisis de Requerimientos:} Identificación de los tipos de datos tensoriales en 
    YOLOv2 y las necesidades de interpolación espacial y visualización volumétrica en Blender.
    \item \textbf{Fase de Diseño Arquitectónico:} Definición del protocolo de intercambio desacoplado 
    (manifiesto JSON y volúmenes OpenVDB) y la separación de responsabilidades entre el componente de 
    extracción y el de visualización.
    \item \textbf{Fase de Desarrollo Iterativo (Prototipado):} 
    \begin{itemize}
        \item \textit{Iteración 1:} Construcción de los scripts de extracción (\textit{forward hooks}) en 
        PyTorch y exportación a matrices numéricas.
        \item \textit{Iteración 2:} Implementación del convertidor a OpenVDB y diseño de la estructura del 
        manifiesto JSON.
        \item \textit{Iteración 3:} Desarrollo del \textit{add-on} en Blender (API \textit{bpy}) para la 
        lectura del manifiesto y la instanciación de nodos.
        \item \textit{Iteración 4:} Diseño de los materiales volumétricos y el sistema de rebanado 
        (\textit{slicing}) en \textit{Geometry Nodes}.
    \end{itemize}
    \item \textbf{Fase de Integración y Pruebas:} Verificación de la consistencia de los datos importados 
    frente a los tensores originales y optimización de rendimiento en la GPU.
\end{enumerate}

\subsection{Estrategia de prueba y validación}
Dado que el propósito central de la investigación es demostrar la **viabilidad técnica del flujo de trabajo 
de visualización** y no evaluar el desempeño del modelo predictivo en producción, la validación del sistema 
se realiza mediante un \textbf{estudio de caso técnico}:

\begin{itemize}
    \item \textbf{Sujeto de prueba:} La arquitectura convolucional YOLOv2 procesando imágenes fijas de prueba.
    \item \textbf{Criterios de evaluación técnica:} 
    \begin{enumerate}
        \item \textit{Fidelidad de los datos:} Comprobación de que las dimensiones, la topología y la densidad 
        escalar representadas en el volumen 3D coincidan numéricamente con los mapas de activación extraídos en 
        PyTorch.
        \item \textit{Resiliencia y robustez:} Verificación de la capacidad del complemento para manejar errores 
        de archivos o capas faltantes degradando la representación de forma segura.
        \item \textit{Interactividad y control visual:} Evaluación cualitativa de la capacidad de respuesta de 
        los nodos de Blender para manipular el mapa de color, la opacidad y los cortes transversales del tensor 
        en tiempo real.
    \end{enumerate}
\end{itemize}